\documentclass[]{manual}

\title{Interference}
\date{\today}
\course{Physics 325}

\begin{document}
    \maketitle
    \thispagestyle{firstpage}

    %%%

    \section*{Motivation}
    
    One of the primary concerns in optics is the measurement and manipulation of the properties of light itself.
    Even other experiments that simply use light to perform some other task typically need to begin by creating that light in the first place.
    The best (and one of the only) ways of measuring the microscopic properties of light - wavelength and frequency - is to exploit that fact that light can \emph{interfere} with itself.


    \section*{Background}
    
    The \emph{superposition principle} of electromagnetism tells us that electric fields can be added to each other as long as they both obey Maxwell's equations independently.
    So, if we have two one-dimensional electric fields, $E_1(x, t)$ and $E_2(x, t)$, the total electric field would be
    \begin{align} \label{eqn:electric-field-superposition}
        E(x, t) = E_1(x, t) + E_2(x, t).
    \end{align}
    Suppose these electric fields belong to a propagating electromagnetic wave (i.e., light) and are shone on a screen.
    What do we expect to see on the screen?
    
    What we see depends on the \emph{intensity} of the light, which is proportional to the square of the electric field:
    \begin{align} \label{eqn:electric-field-intensity}
        I(x, t) &= \epsilon_0 c \, E^2(x, t) \nonumber \\
        I(x, t) &= \epsilon_0 c \,  \qty[ E_1^2(x, t) + E_2^2(x, t) + E_1(x, t) E_2(x, t) ]
    \end{align}
    Equation \eqref{eqn:electric-field-intensity} is appropriate in a vacuum, where $\epsilon_0$ is the vacuum permittivity and $c$ is the speed of light.
    The first two terms are what we would expect from the intensity of the individual fields.
    The last term is the interference: it depends on the product of the electric fields, so it could be positive or negative or zero at any given point.
    

    \section*{Equipment}


    \newpage

    \section*{Procedure}
    
    \subsection{Collimating the Beam}
    For Fraunhofer diffraction, we assumed that plane waves were incident on our slits.
    To get as close to a plane wave as possible, you must use the telescope to collimate the He-Ne beam.
    Place the telescope in front of the laser (make sure you put the correct end of the telescope toward the laser).
    Adjust the height and orientation of the telescope until you can see a circular beam spot in the center of the screen (about \SI{1}{\meter} down the optical rail).
    Adjust the focus of the telescope until the beam is approximately the same diameter for the entire distance between the telescope and the screen (use a handy piece of paper to look at the beam at any point in space).
    You might also need to change the orientation of the laser to get the beam centered on the screen.
    


    \newpage
    \section*{Prelab Questions}
    \setcounter{subsection}{0}

    \subsection{Question 1 \hfill (1 pt.)}

    \vspace{\stretch{1}}

    \subsection{Question 2 \hfill (1 pt.)}

    \vspace{\stretch{1}}

\end{document}
